\chapter*{Index of Examples}
\addcontentsline{toc}{chapter}{Index of Examples}
\phantomsection

\noindent Examples are numbered as \textbf{Chapter.Lecture.Example}.

\begin{multicols}{2}
\small

%%%%%%%%%%%%%%%%%%%%%%%%%%%%%%%%%%%%%%%%%%%%%%%%%%
\subsection*{Chapter 1}
%%%%%%%%%%%%%%%%%%%%%%%%%%%%%%%%%%%%%%%%%%%%%%%%%%

\noindent\textit{Lecture 1.1 --- Review of Differential and Integral Calculus}\\
\textbf{1.1.1} Product rule for $x^2 e^x$ \(\bullet\)
\textbf{1.1.2} Integral of $1/(1-x)$ \(\bullet\)
\textbf{1.1.3} Solving a system of equations \(\bullet\)
\textbf{1.1.4} L'H\^opital's Rule evaluation \(\bullet\)
\textbf{1.1.5} Epsilon-delta proof of a limit \(\bullet\)
\textbf{1.1.6} Limit definition of the derivative \(\bullet\)
\textbf{1.1.7} Simpson's Rule approximation \(\bullet\)
\textbf{1.1.8} Integration by DI (Tabular) Method \(\bullet\)
\textbf{1.1.9} Taylor series of $\ln(1+x)$ \(\bullet\)
\textbf{1.1.10} Mean Value Theorem verification

\medskip
\noindent\textit{Lecture 1.2 --- Three-Dimensional Coordinate Systems}\\
\textbf{1.2.1} Distance between two points in 3D \(\bullet\)
\textbf{1.2.2} Cylindrical to rectangular conversion \(\bullet\)
\textbf{1.2.3} Spherical to rectangular conversion

\medskip
\noindent\textit{Lecture 1.3 --- Vectors}\\
\textbf{1.3.1} Magnitude, unit vector, direction cosines \(\bullet\)
\textbf{1.3.2} Adding and scaling force vectors \(\bullet\)
\textbf{1.3.3} Dot product of two vectors \(\bullet\)
\textbf{1.3.4} Cross product computations \(\bullet\)
\textbf{1.3.5} Cross product of two given vectors \(\bullet\)
\textbf{1.3.6} Cosine similarity of feature vectors

\medskip
\noindent\textit{Lecture 1.4 --- Lines and Planes in Space}\\
\textbf{1.4.1} Parametric equations of a line \(\bullet\)
\textbf{1.4.2} Equation of a plane from point and normal \(\bullet\)
\textbf{1.4.3} Line--plane intersection \(\bullet\)
\textbf{1.4.4} Intersection of two planes \(\bullet\)
\textbf{1.4.5} Projection onto a plane \(\bullet\)
\textbf{1.4.6} Plane through three points \(\bullet\)
\textbf{1.4.7} Distance from a point to a plane

\medskip
\noindent\textit{Lecture 1.5 --- Quadric Surfaces}\\
\textbf{1.5.1} Plane intersecting a quadric surface \(\bullet\)
\textbf{1.5.2} Classifying a quadric by its equation \(\bullet\)
\textbf{1.5.3} Cylinder--plane intersection

%%%%%%%%%%%%%%%%%%%%%%%%%%%%%%%%%%%%%%%%%%%%%%%%%%
\subsection*{Chapter 2: Vector-Valued Functions and Motion in Space}
%%%%%%%%%%%%%%%%%%%%%%%%%%%%%%%%%%%%%%%%%%%%%%%%%%

\noindent\textit{Lecture 2.1 --- Vector-Valued Functions}\\
\textbf{2.1.1} Sketching a vector-valued curve \(\bullet\)
\textbf{2.1.2} Domain of a vector-valued function \(\bullet\)
\textbf{2.1.3} Parametric curve analysis

\medskip
\noindent\textit{Lecture 2.2 --- Motion Along a Curve}\\
\textbf{2.2.1} Position from a velocity vector \(\bullet\)
\textbf{2.2.2} Projectile motion over a sandtrap \(\bullet\)
\textbf{2.2.3} 3D projectile with given launch angle

\medskip
\noindent\textit{Lecture 2.3 --- The TNB Frame, Curvature and Torsion}\\
\textbf{2.3.1} TNB frame for a space curve \(\bullet\)
\textbf{2.3.2} Curvature of a curved road design \(\bullet\)
\textbf{2.3.3} Torsion of a space curve

\medskip
\noindent\textit{Lecture 2.4 --- Arc Length and Reparameterization}\\
\textbf{2.4.1} Arc length of a circular sector \(\bullet\)
\textbf{2.4.2} Arc length of a parametric curve \(\bullet\)
\textbf{2.4.3} Arc length and reparameterization of a helix \(\bullet\)
\textbf{2.4.4} Arc length of a nearly flat curve

%%%%%%%%%%%%%%%%%%%%%%%%%%%%%%%%%%%%%%%%%%%%%%%%%%
\subsection*{Chapter 3: Partial Derivatives}
%%%%%%%%%%%%%%%%%%%%%%%%%%%%%%%%%%%%%%%%%%%%%%%%%%

\noindent\textit{Lecture 3.1 --- Functions of Several Variables}\\
\textbf{3.1.1} Writing famous two-variable equations \(\bullet\)
\textbf{3.1.2} Domain and range of multivariable functions \(\bullet\)
\textbf{3.1.3} Level curves of $\sin(x^2+y^2)$ \(\bullet\)
\textbf{3.1.4} Level curves of $\cos(y)$

\medskip
\noindent\textit{Lecture 3.2 --- Limits and Continuity}\\
\textbf{3.2.1} Evaluating multivariable limits \(\bullet\)
\textbf{3.2.2} Limits via advanced techniques \(\bullet\)
\textbf{3.2.3} Epsilon-delta limit proof in $\mathbb{R}^2$ \(\bullet\)
\textbf{3.2.4} Intermediate Value Theorem verification

\medskip
\noindent\textit{Lecture 3.3 --- Partial Derivatives}\\
\textbf{3.3.1} Partial derivative via limit definition \(\bullet\)
\textbf{3.3.2} Computing partial derivatives \(\bullet\)
\textbf{3.3.3} Verifying Clairaut's Theorem \(\bullet\)
\textbf{3.3.4} Implicit partial derivatives for a container

\medskip
\noindent\textit{Lecture 3.4 --- The Chain Rule}\\
\textbf{3.4.1} Multivariable Chain Rule computations \(\bullet\)
\textbf{3.4.2} Second-order mixed partial via Chain Rule \(\bullet\)
\textbf{3.4.3} Implicit differentiation in three variables \(\bullet\)
\textbf{3.4.4} Total differential for temperature change

\medskip
\noindent\textit{Lecture 3.5 --- Directional Derivatives and Gradient Vectors}\\
\textbf{3.5.1} Computing a gradient vector \(\bullet\)
\textbf{3.5.2} Directional derivative verification \(\bullet\)
\textbf{3.5.3} Steepest descent via level curves \(\bullet\)
\textbf{3.5.4} Gradient application to temperature field

\medskip
\noindent\textit{Lecture 3.6 --- Tangent Planes and Linear Approximations}\\
\textbf{3.6.1} Tangent plane to a surface \(\bullet\)
\textbf{3.6.2} Linear approximation for terrain elevation

\medskip
\noindent\textit{Lecture 3.7 --- Extrema and Critical Points}\\
\textbf{3.7.1} Extreme Value Theorem illustration \(\bullet\)
\textbf{3.7.2} Critical points via First Derivative Test \(\bullet\)
\textbf{3.7.3} Second Derivative Test at a critical point \(\bullet\)
\textbf{3.7.4} Global extrema on a closed region \(\bullet\)
\textbf{3.7.5} Saddle point of a product of sines

\medskip
\noindent\textit{Lecture 3.8 --- Optimization}\\
\textbf{3.8.1} Minimizing a manufacturing cost function \(\bullet\)
\textbf{3.8.2} Maximum volume box under a plane \(\bullet\)
\textbf{3.8.3} Lagrange multiplier with one constraint \(\bullet\)
\textbf{3.8.4} Constrained production cost minimization \(\bullet\)
\textbf{3.8.5} Dietary risk minimization \(\bullet\)
\textbf{3.8.6} Support Vector Machine classification

%%%%%%%%%%%%%%%%%%%%%%%%%%%%%%%%%%%%%%%%%%%%%%%%%%
\subsection*{Chapter 4: Multiple Integrals}
%%%%%%%%%%%%%%%%%%%%%%%%%%%%%%%%%%%%%%%%%%%%%%%%%%

\noindent\textit{Lecture 4.1 --- Double Integrals over Rectangular Regions}\\
\textbf{4.1.1} Riemann sum approximation of a double integral \(\bullet\)
\textbf{4.1.2} Evaluating iterated double integrals \(\bullet\)
\textbf{4.1.3} Comparing integration orders \(\bullet\)
\textbf{4.1.4} Indefinite double integrals

\medskip
\noindent\textit{Lecture 4.2 --- Double Integrals over General Regions}\\
\textbf{4.2.1} Finding bounds for a curved region \(\bullet\)
\textbf{4.2.2} Average efficiency of a solar panel \(\bullet\)
\textbf{4.2.3} Potential flux through a curved channel \(\bullet\)
\textbf{4.2.4} Reversing integration order for $\sin(y^2)$ \(\bullet\)
\textbf{4.2.5} Double integral in polar coordinates \(\bullet\)
\textbf{4.2.6} Basin excavation with angular depth \(\bullet\)
\textbf{4.2.7} The Gaussian integral \(\bullet\)
\textbf{4.2.8} Area of a region via double integral

\medskip
\noindent\textit{Lecture 4.3 --- Applications of Double Integrals}\\
\textbf{4.3.1} Center of mass of a solar sail \(\bullet\)
\textbf{4.3.2} Turntable platter with variable density \(\bullet\)
\textbf{4.3.3} Moment of inertia of a mounting plate \(\bullet\)
\textbf{4.3.4} Surface area of a hyperbolic paraboloid

\medskip
\noindent\textit{Lecture 4.4 --- Triple Integrals}\\
\textbf{4.4.1} Triple integral over a rectangular box \(\bullet\)
\textbf{4.4.2} Volume of intersecting cylinders \(\bullet\)
\textbf{4.4.3} Volume of a cone in cylindrical coordinates \(\bullet\)
\textbf{4.4.4} Cylindrical beam inside a spherical dome

\medskip
\noindent\textit{Lecture 4.5 --- Change of Variables in Multiple Integrals}\\
\textbf{4.5.1} Jacobian for toroidal coordinates \(\bullet\)
\textbf{4.5.2} Linear transformation over a parallelogram \(\bullet\)
\textbf{4.5.3} Spiral galaxy density integration

\medskip
\noindent\textit{Lecture 4.6 --- Applications of Triple Integrals}\\
\textbf{4.6.1} Center of mass of a cylindrical rod \(\bullet\)
\textbf{4.6.2} Lorentz force on a conductive fluid \(\bullet\)
\textbf{4.6.3} Total force on an expanding gas volume \(\bullet\)
\textbf{4.6.4} Average temperature in a blast furnace

%%%%%%%%%%%%%%%%%%%%%%%%%%%%%%%%%%%%%%%%%%%%%%%%%%
\subsection*{Chapter 5: Vector Calculus}
%%%%%%%%%%%%%%%%%%%%%%%%%%%%%%%%%%%%%%%%%%%%%%%%%%

\noindent\textit{Lecture 5.1 --- Vector Fields}\\
\textbf{5.1.1} Classifying and visualizing vector fields \(\bullet\)
\textbf{5.1.2} Conservative vs.\ non-conservative fields

\medskip
\noindent\textit{Lecture 5.2 --- Line Integrals}\\
\textbf{5.2.1} Scalar line integral along two paths \(\bullet\)
\textbf{5.2.2} Scalar function line integral \(\bullet\)
\textbf{5.2.3} Vector line integral on a circle \(\bullet\)
\textbf{5.2.4} Piecewise path work integral \(\bullet\)
\textbf{5.2.5} Testing for conservative vector fields \(\bullet\)
\textbf{5.2.6} Finding a potential function

\medskip
\noindent\textit{Lecture 5.3 --- Green's Theorem}\\
\textbf{5.3.1} Circulation via Green's Theorem \(\bullet\)
\textbf{5.3.2} Vorticity integral over a basin

\medskip
\noindent\textit{Lecture 5.4 --- Curl and Divergence}\\
\textbf{5.4.1} Curl of an oscillatory electric field \(\bullet\)
\textbf{5.4.2} Conservative fields are irrotational \(\bullet\)
\textbf{5.4.3} Divergence of a pipe flow field \(\bullet\)
\textbf{5.4.4} Vector operator identities \(\bullet\)
\textbf{5.4.5} Curl in cylindrical coordinates

\medskip
\noindent\textit{Lecture 5.5 --- Surface Integrals}\\
\textbf{5.5.1} Parametric and scalar surface conversions \(\bullet\)
\textbf{5.5.2} Techniques for surface representations \(\bullet\)
\textbf{5.5.3} Surface area of a paraboloid \(\bullet\)
\textbf{5.5.4} Average efficiency of a curved solar panel \(\bullet\)
\textbf{5.5.5} Mass of coating on a hemispherical dome \(\bullet\)
\textbf{5.5.6} Flux through a cylindrical air filter

\medskip
\noindent\textit{Lecture 5.6 --- Stokes' Theorem}\\
\textbf{5.6.1} Stokes' Theorem on a paraboloid \(\bullet\)
\textbf{5.6.2} Line integral on a plane--cylinder curve \(\bullet\)
\textbf{5.6.3} Divergence Theorem on a hemisphere \(\bullet\)
\textbf{5.6.4} Amp\`ere's law via Stokes' Theorem \(\bullet\)
\textbf{5.6.5} Stokes' Theorem on a unit circle

\medskip
\noindent\textit{Lecture 5.7 --- Divergence Theorem}\\
\textbf{5.7.1} Divergence Theorem on a paraboloid region \(\bullet\)
\textbf{5.7.2} Gravitational field of a variable-density sphere

%%%%%%%%%%%%%%%%%%%%%%%%%%%%%%%%%%%%%%%%%%%%%%%%%%
\subsection*{Chapter 6: Differential Equations}
%%%%%%%%%%%%%%%%%%%%%%%%%%%%%%%%%%%%%%%%%%%%%%%%%%

\noindent\textit{Lecture 6.1 --- Introduction to Differential Equations}\\
\textbf{6.1.1} Newton's Law of Cooling for coffee \(\bullet\)
\textbf{6.1.2} Solving $dy/dx = -2xy$ by separation \(\bullet\)
\textbf{6.1.3} Verifying a series solution to an ODE

\medskip
\noindent\textit{Lecture 6.2 --- First-Order Ordinary Differential Equations}\\
\textbf{6.2.1} Separable ODE with factoring \(\bullet\)
\textbf{6.2.2} First-order linear ODE via integrating factor \(\bullet\)
\textbf{6.2.3} Exact differential equation \(\bullet\)
\textbf{6.2.4} Bernoulli equation $y' + y = y^2$

\medskip
\noindent\textit{Lecture 6.3 --- Higher-Order ODEs I}\\
\textbf{6.3.1} Second-order non-homogeneous ODE \(\bullet\)
\textbf{6.3.2} Fifth-order homogeneous ODE

\medskip
\noindent\textit{Lecture 6.4 --- Higher-Order ODEs II}\\
\textbf{6.4.1} Variation of Parameters

\medskip
\noindent\textit{Lecture 6.5 --- Higher-Order ODEs III}\\
\textbf{6.5.1} Inverse Laplace transform \(\bullet\)
\textbf{6.5.2} IVP via Laplace transforms (sine forcing) \(\bullet\)
\textbf{6.5.3} IVP via Laplace transforms (delta forcing)

\medskip
\noindent\textit{Lecture 6.6 --- Methods of Solving PDEs I}\\
\textbf{6.6.1} Heat equation with Dirichlet conditions \(\bullet\)
\textbf{6.6.2} Vibrating membrane on a circular drum \(\bullet\)
\textbf{6.6.3} Wave equation via separation of variables \(\bullet\)
\textbf{6.6.4} Poisson's equation on a rectangle

\medskip
\noindent\textit{Lecture 6.7 --- Methods of Solving PDEs II}\\
\textbf{6.7.1} PDE via method of characteristics \(\bullet\)
\textbf{6.7.2} First-order PDE with initial condition \(\bullet\)
\textbf{6.7.3} 3D Poisson equation with Gaussian source \(\bullet\)
\textbf{6.7.4} Schr\"odinger equation Green's function

%%%%%%%%%%%%%%%%%%%%%%%%%%%%%%%%%%%%%%%%%%%%%%%%%%
\subsection*{Chapter 7: Complex Variables}
%%%%%%%%%%%%%%%%%%%%%%%%%%%%%%%%%%%%%%%%%%%%%%%%%%

\noindent\textit{Lecture 7.1 --- The Complex Number System}\\
\textbf{7.1.1} Complex fraction in standard form \(\bullet\)
\textbf{7.1.2} Geometric locus of complex equations \(\bullet\)
\textbf{7.1.3} Computing $i^i$ \(\bullet\)
\textbf{7.1.4} Cube roots of $-8$ in the complex plane \(\bullet\)
\textbf{7.1.5} Primitive roots of $z^4 - 1 = 0$

\medskip
\noindent\textit{Lecture 7.2 --- Complex Functions}\\
\textbf{7.2.1} Poles and orders of a rational function \(\bullet\)
\textbf{7.2.2} Solving $e^z = 1 + i$ \(\bullet\)
\textbf{7.2.3} Solving $\cos(z) = 2$ \(\bullet\)
\textbf{7.2.4} All values of $\ln(i)$ \(\bullet\)
\textbf{7.2.5} Principal value of $(-8)^{1/3}$

\medskip
\noindent\textit{Lecture 7.3 --- Basic Limits in Complex Analysis}\\
\textbf{7.3.1} Evaluating complex limits \(\bullet\)
\textbf{7.3.2} Continuity of $|z|^2$ \(\bullet\)
\textbf{7.3.3} Open mapping of $z^2$ \(\bullet\)
\textbf{7.3.4} Connectedness of $e^{\mathbb{C}}$ \(\bullet\)
\textbf{7.3.5} Compactness of a closed annulus \(\bullet\)
\textbf{7.3.6} Limit of $1/z$ in extended plane \(\bullet\)
\textbf{7.3.7} Sequential discontinuity of $\bar{z}/z$

\medskip
\noindent\textit{Lecture 7.4 --- Complex Differentiation}\\
\textbf{7.4.1} Derivative via limit definition \(\bullet\)
\textbf{7.4.2} Cauchy--Riemann equations verification \(\bullet\)
\textbf{7.4.3} Complex differentiation rules \(\bullet\)
\textbf{7.4.4} Derivative of $\mathrm{Log}(z)$ \(\bullet\)
\textbf{7.4.5} Laurent series of $1/[z(1-z)]$ \(\bullet\)
\textbf{7.4.6} Liouville's Theorem application

\medskip
\noindent\textit{Lecture 7.5 --- Elementary Complex Integration}\\
\textbf{7.5.1} Line integral of $z^2$ on a segment \(\bullet\)
\textbf{7.5.2} Integral of $\bar{z}$ on the unit circle \(\bullet\)
\textbf{7.5.3} Contour integral and Jordan's Lemma \(\bullet\)
\textbf{7.5.4} Winding number integral of $1/z$ \(\bullet\)
\textbf{7.5.5} Antiderivative and path independence

\medskip
\noindent\textit{Lecture 7.6 --- Cauchy's Integral Theorem and Residues}\\
\textbf{7.6.1} Cauchy's Theorem with the Gamma function \(\bullet\)
\textbf{7.6.2} Non-holomorphicity via Morera's Theorem \(\bullet\)
\textbf{7.6.3} Cauchy's Integral Formula application \(\bullet\)
\textbf{7.6.4} Higher-order Cauchy formula evaluation \(\bullet\)
\textbf{7.6.5} Laurent series and singularity classification \(\bullet\)
\textbf{7.6.6} Residue computation with multiple poles \(\bullet\)
\textbf{7.6.7} Essential singularity residue integral \(\bullet\)
\textbf{7.6.8} Trigonometric integral via residues
\end{multicols}
